%\documentclass[a4paper,12pt,answers]{exam} % énoncé+réponses penser à commenter le multicol
\documentclass[a4paper,10pt]{book}        % énoncé seul
\include{../1ere-preambule}

\begin{document}

%\setcounter{chapter}{}
\vspace{-0.6 cm}
\chapter{Automatismes - Exercices}
\thispagestyle{fancy}

\vspace{-0.8 cm}
\begin{minipage}[c]{0.60\linewidth}
\exercice{(Sous forme décimale ou sous forme de pourcentage)}
\begin{enumMet}
	\item Écrire sous forme décimale :
	\begin{enumMet}
		\begin{multicols}{5}
			\item $7,5 \%$ ;
		\item $0,75\%$ ;
		\item $0,05\%$ ;
		\item $50\%$ ;
		\item $200\%$.
		\end{multicols}
	\end{enumMet}
	\item Écrire sous forme de pourcentage:
	\begin{enumMet}
		\begin{multicols}{5}
			\item $0,40$;
		\item $0,08$;
		\item $0,015$;
		\item $0,0075$;
		\item $0,001$.
		\end{multicols}
	\end{enumMet}
\end{enumMet}
	
\end{minipage} \hfill
\begin{minipage}[c]{0.38\linewidth}
	\exercice{(Calculer un pourcentage)}\\
	Calculer:
	%\setlength{\columnseprule}{0pt}
%	\begin{multicols}{2}
		\begin{enumMet}
			\item $25 \%$ de $300$ ;
			\item $33 \%$ de $700$;
			\item $0,5 \%$ de $2 496 000$;
			\item $300 \%$ de $12$.
		\end{enumMet}
%	\end{multicols}
\end{minipage}

\ \\
\mbox{}\hfill\rule{0.8\linewidth}{1pt}\hfill\mbox{}

\begin{minipage}[c]{0.58\linewidth}
\exercice{(Calculer un pourcentage)}\\
Calculer:
\begin{multicols}{2}
	\begin{enumMet}
		\item $10 \%$ de $5 000$ ;
		\item $20 \%$ de $6,50$ ;
		\item $200 \%$ de $2$;
		\item $0,5 \%$ de $50 000$;
	\end{enumMet}
\end{multicols}	

\exercice{(Un rabais sur l'assurance automobile)} \\
Pour un certain modèle de voiture, un assureur propose $30\%$
d'économie moyenne par rapport à l'assureur précédent, soit
$150$ \EUR{} d'économie. Quel tarif cet assureur attribue-t-il, pour
ce modèle de voiture, à l'assureur précédent?
\end{minipage} \hfill
\begin{minipage}[c]{0.38\linewidth}
\exercice{ $-$ Salle informatique}

Deux salles informatiques du lycée sont équipées de la façon suivante:

\begin{itemMet}
	\item la salle 1 avec 25 ordinateurs, dont 6 neufs ;
	\item la salle 2 avec 20 ordinateurs, dont 5 neufs.
\end{itemMet}
	\guillemets{Je préfère aller en salle 1, c'est là qu'il y a le plus d'ordinateurs neufs} dit une élève.
	
	\guillemets{La proportion d'ordinateurs neufs est plus grande dans la salle 2} lui répond son professeur. Qui a raison ?
\end{minipage}


\ \\
\mbox{}\hfill\rule{0.8\linewidth}{1pt}\hfill\mbox{}


\begin{minipage}[c]{0.48\linewidth}
	\exercice{(Calculer une proportion ou un effectif)}
	\begin{enumMet}
		\item Calculer $p$ lorsque $n_S = 150$ et $n_E = 600$.\\
		Écrire le résultat sous la forme d'un pourcentage.
		\item Calculer $p$ lorsque $n_S = 12$ et $n_E= 2 400$.\\
		Écrire le résultat sous la forme d'un pourcentage.
		\item Calculer $n_S$ lorsque $p = 0,09$ et $n_E= 250 000$.
		\item Calculer $n_E$ lorsque $p = 0,30$ et $n_S = 3 000$.
		\item Calculer $n_E$ lorsque $p = 20 \%$ et $n_S = 6 000$.
		\item Sachant que $30 \%$ d'une somme $S$ vaut $600$ \EUR{}, calculer $S$.
	\end{enumMet}
\end{minipage} \hfill
\begin{minipage}[c]{0.48\linewidth}
	\exercice{(Après le bac)}
	
	Dans un lycée technologique, $200$ élèves se sont présentés
	au baccalauréat et $90 \%$ des élèves ont été reçus.\\
	\begin{enumMet}
		\item Calculer le nombre d'élèves reçus.
		\item Parmi les élèves reçus $20 \%$ se sont inscrits en IUT, cinq sont partis dans la vie active, les autres ont été admis en sections de technicien supérieur.
		\begin{enumMet}
			\item Calculer le nombre d'élèves inscrits dans un IUT.\\
			\item Calculer le nombre d'élèves inscrits dans une section de technicien supérieur.
		\end{enumMet}
	\end{enumMet}
\end{minipage}

\ \\
\mbox{}\hfill\rule{0.8\linewidth}{1pt}\hfill\mbox{}

\begin{minipage}[c]{0.38\linewidth}
	\exercice{ $-$ Taux d'évolution entre $Q_{1}$ et $Q_{2}$}
	\begin{enumMet}
		\item Calculer l'un des trois nombres $Q_{1}, Q_{2}$ ou $t$, connaissant les deux autres.
		\item Indiquer s'il s'agit d'une hausse ou d'une baisse et donner le taux d'évolution sous forme de pourcentage.
		\begin{enumMet}
			\item $Q_{1}=3 ; Q_{2}=4$.
			\item $Q_{1}=2,5 ; Q_{2}=2$.
			\item $Q_{1}=3 ; t=-0,10$.
			\item $Q_{2}=1,2 ; t=0,2$.
		\end{enumMet}
	\end{enumMet}
\end{minipage} \hfill
\begin{minipage}[c]{0.68\linewidth}
	\exercice{ $-$ Pourcentage et coef multiplicateur}
	\begin{enumMet}
		\item Dans chacun des cas suivants, donner le coefficient multiplicateur correspondant à une hausse ou à une baisse de pourcentage donné.
		\begin{multicols}{2}
			\begin{enumMet}
			\item une hausse de $30 \%$;
			\item une baisse de $30 \%$;
			\item une hausse de $45 \%$;
			\item une baisse de $91 \%$;
			\item une hausse de $300 \%$.
		\end{enumMet}
		\end{multicols}
		\item Dans chacun des cas suivants, le coefficient multiplicateur $c$ est donné. Indiquer sil s'agit d'une hausse ou d'une baisse et en donner le pourcentage.
		\begin{multicols}{2}
			\begin{enumMet}
			\item $c=1,03$;
			\item $c=1,025$;
			\item $c=0,2$;
			\item $c=0,70$;
			\item $c=2,5$;
			\item $c=0,995$.
		\end{enumMet}
		\end{multicols}
	\end{enumMet}
\end{minipage}

\newpage

\begin{minipage}[c]{0.38\linewidth}
	\exercice{}
	
	Recopier et compléter les phrases suivantes.
	\begin{enumMet}
		\item Augmenter de $40 \%$, c'est multiplier par...
		\item Diminuer de $50 \%$, c'est multiplier par...
		\item Multiplier par 3 c'est augmenter de ...\%.
		\item Multiplier par 0,25 c'est diminuer de ...\%.	
	\end{enumMet}
\end{minipage} \hfill
\begin{minipage}[c]{0.58\linewidth}
	\exercice{}
	
	\begin{enumMet}
		\item Une personne paie, pour un groupe, une note de restaurant qui s'élève à 207 \EUR{}, avec le service compris de $15 \%$. Quel est le prix des repas sans le service?
		\item Un commerçant calcule ses prix de vente en prenant un bénéfice de $30 \%$ sur ses prix d'achat. Quel est le prix d'achat d'un article qu'il a vendu $113,10$\EUR{}.
		\item Le prix d'un article soldé est de $41,40$ \EUR{}. L'étiquette indique \guillemets{$-40 \%$}. Calculer le prix de l'article avant les soldes.
	\end{enumMet}
\end{minipage}

\ \\
\mbox{}\hfill\rule{0.8\linewidth}{1pt}\hfill\mbox{}

\begin{minipage}[c]{0.48\linewidth}
	\exercice{}
	
	Dans chacun des cas suivants, calculer le coefficient multiplicateur global. Indiquer s'il s'agit d'une baisse ou d'une hausse et en donner l'e taux d'évolution sous forme de pourcentage.
	\begin{enumMet}
		\item une hausse de $10 \%$, puis une baisse de $20 \%$;
		\item une hausse de $20 \%$, puis une baisse de $10 \%$;
		\item une hausse de $10 \%$, puis une hausse de $10 \%$;
		\item une baisse de $10 \%$, puis une baisse de $10 \%$;
		\item une baisse de $50 \%$, puis une hausse de $200 \%$.
	\end{enumMet}
\end{minipage} \hfill
\begin{minipage}[c]{0.48\linewidth}
	\exercice{}
	
	\begin{enumMet}
		\item Un prix augmente de $40 \%$, puis baisse de $40 \%$. Le prix est-il revenu à sa valeur initiale ?
		\item Deux hausses successives de $50 \%$ sont-elles équivalentes à une hausse de $100 \%$ ?
		\item Une baisse de $10 \%$ suivie d'une baisse de $20 \%$ a-t-elle le même effet qu'une baisse de $20 \%$ suivie d'une baisse de $10 \%$ ?
	\end{enumMet}
\end{minipage}

\ \\
\mbox{}\hfill\rule{0.8\linewidth}{1pt}\hfill\mbox{}


\exercice{}

Les questions suivantes sont indépendantes.
\begin{enumMet}
	\item Un article qui valait $92 €$ a subi deux augmentations successives, la première de $5 \%$, la seconde de $15 \%$.\\
	Quelle est l'augmentation totale, en pourcentage et en valeur, subie par cet article?
	\item Après deux augmentations successives, la première de $10 \%$, la seconde de $20 \%$, un matériel coûte $792 €$.\\
	Combien coûtait-il avant les deux augmentations?
	\item Le prix d'un produit industriel a d'abord augmenté de $8,5 \%$ puis diminué de $3 \%$. Calculer le taux d'évolution global de ce prix sous forme de pourcentage. Arrondir à $0,01 \%$.
	\item Le prix d'un produit alimentaire a subi trois hausses mensuelles successives de $9,5 \%$. Calculer le taux d'évolution global sous forme de pourcentage. Arrondir à $0,1 \%$.
	\item La valeur de revente d'un équipement informatique pour les entreprises a subi quatre baisses annuelles successives de $25 \%$. Au bout de quatre ans la valeur de revente de l'équipement:
	\begin{multicols}{2}
		\begin{enumMet}
		\item est nulle;
		\item a baissé d'environ $60 \%$;
		\item a baissé d'environ $68 \%$;
		\item a baissé d'environ $75 \%$.
	\end{enumMet}
	\end{multicols}
	
\end{enumMet}

\ \\
\mbox{}\hfill\rule{0.8\linewidth}{1pt}\hfill\mbox{}


\exercice{ $-$ Le prix d'une matière première}
Une matière première coûtait $140$ \EUR{} la tonne au début du mois d'août dernier.

\begin{enumMet}
	\item Calculer le prix $P_{1}$ à la tonne de cette matière première après une hausse de $7 \%$.
	\item \begin{enumMet}
		\item Calculer le coefficient multiplicateur $c_{1}$ de la baisse qu'il faudrait appliquer au prix de la tonne de cette matière première pour qu'il revienne à $140$ \EUR{}. Arrondir à $10^{-4}$.
		\item En déduire le taux d'évolution $t_{1}$ de cette baisse sous forme de pourcentage.
	\end{enumMet}
\end{enumMet}

\ \\
\mbox{}\hfill\rule{0.8\linewidth}{1pt}\hfill\mbox{}


\exercice{ $-$ Évolutions successives et évolution réciproque}
\begin{enumMet}
	\item Pendant une période de promotion, le prix d'un outillage a subi deux baisses successives de $12 \%$. Démontrer que le taux d'évolution de la hausse, qu'il faut appliquer au prix de l'outillage pour revenir au prix initial avant la période de promotion, est, en arrondissant à $0,01 \%$, de $29,13 \%$.
	\item Le prix d'un article a subi deux hausses successives de $12 \%$. Déterminer le taux d'évolution de la baisse qu'il faut appliquer au prix de l'article pour revenir au prix initial.
\end{enumMet}

\newpage

\exercice{ $-$ Somme de fractions}

\begin{multicols}{2}
	\begin{enumMet}
	\item Trouver une fraction irréductible égale à $\dfrac{4}{7}+\dfrac{3}{8}$.
	\item Trouver une fraction irréductible égale à $3-\dfrac{2}{3}$.
	\item Vérifier que : $\dfrac{1}{3} x+\dfrac{2}{5}=\dfrac{5 x+6}{15}$.
	\item Réduire: $A=\dfrac{4}{3} x-2 x$.
	\item Reproduire et compléter l'égalité suivante : $0,72=1-\dfrac{\cdots}{100}$.
\end{enumMet}
\end{multicols}

\exercice{ $-$ Produit de fractions}
\begin{multicols}{2}
	\begin{enumMet}
	\item Calculer: $p=\dfrac{5}{12} \times \dfrac{16}{15}$. Simplifier le résultat obtenu.
	\item Trouver une fraction irréductible égale à $\dfrac{15}{4} \times \dfrac{8}{25}$.
	\item Calculer : $p=90 \times \dfrac{30}{100}$.
	\item Calculer : $p=40 \times\left(1+\dfrac{10}{100}\right)$.
	\item Développer : $p=\dfrac{1}{2}\left(\dfrac{3}{4} x+\dfrac{1}{3}\right)$.\\
\end{enumMet}
\end{multicols}

\exercice{ $-$ notation scientifique}

Écrire en notation scientifique les nombres figurant dans chacune des phrases suivantes.
\begin{enumMet}
	\item Le big-bang marquant l'origine de l'univers a eu lieu il y a environ 15 milliards d'années.
	\item La vitesse de la lumière est d'environ $299792000 \mathrm{~m} \cdot \mathrm{~s}^{-1}$.
	\item L'aire de la Terre est d'environ $510100000 \mathrm{~km}^{2}$.
	\item L'étoile la plus proche du soleil est à environ 37500000000000 km .
\end{enumMet}
Écrire sous forme décimale chacun des nombres suivants.
\begin{multicols}{2}
	\begin{enumMet}
	\item $3,42 \times 10^{3}$;
	\item $0,501 \times 10^{4}$;
	\item $17,5 \times 10^{-2}$;
	\item $0,02 \times 10^{-3}$.
\end{enumMet}
\end{multicols}

Écrire sous forme décimale, puis sous forme scientifique, chacun des nombres suivants.
\begin{enumMet}
	\item $\mathrm{A}=3 \times 10^{3}+2 \times 10^{-1}+5 \times 10^{-2}$;
	\item $\mathrm{B}=\dfrac{3 \times 10^{2} \times 4 \times 10^{3}}{2 \times 10^{-3}}$.
\end{enumMet}

\exercice{ $-$ équation du premier degré}

Résoudre une équation du premier degré ou une équation se ramenant à des équations du premier degré * Résoudre dans $\mathbb{R}$ chacune des équations suivantes.
\begin{multicols}{3}
	\begin{enumMet}
	\item $0,2 x=7$;
	\item $12 x=0$;
	\item $4 x=6 x+1$;
	\item $x=0,2 x-1,6$;
	\item $2 x+5=x-2+\dfrac{1}{2} x$;
	\item $\dfrac{3}{4} x=\dfrac{1}{2} x+\dfrac{1}{3}$;
	\item $4(x-6)(x-1)=0$;
	\item $-2(3 x+3)(-x+5)=0$.
\end{enumMet}
\end{multicols}

\exercice{ $-$ Résoudre une inéquation du premier degré}

Résoudre dans \R chacune des inéquations suivantes.
\begin{multicols}{3}
	\begin{enumMet}
	\item $-3 x \geqslant 1$;
	\item $-4 x+7 \leqslant 0$;
	\item $\dfrac{3}{4} x+\dfrac{1}{2} \geqslant 0$;
	\item $-\dfrac{1}{3} x-\dfrac{2}{5} \leqslant 0$;
	\item $6 x-3 \leqslant-2 x+7$;
	\item $-4 x+9 \geqslant \dfrac{1}{3} x+7$.
\end{enumMet}
\end{multicols}

\exercice{ $-$ Quotient de fractions}

\begin{multicols}{2}
	\begin{enumMet}
	\item Calculer : $q=\dfrac{3}{5}\div \dfrac{4}{7}$.
	\item Calculer $: q=3\div\dfrac{4}{5}$.
	\item Calculer : $q=\dfrac{2}{3}\div\dfrac{4}{9}$. Simplifier le résultat obtenu.
	\item Calculer: $q=\dfrac{2-\frac{3}{5}}{4-\frac{7}{5}}$. Mettre le résultat sous forme de fraction irréductible.
	\item Reproduire et compléter l'égalité suivante $: \dfrac{2}{7} \times \ldots=\dfrac{3}{4}$.
\end{enumMet}
\end{multicols}

\exercice{ $-$ Compléter des égalités}

Reproduire et compléter chacune des égalités suivantes.
\begin{multicols}{3}
	\begin{enumMet}
	\item $10 \cdots \times 10^{4}=10^{-2}$;
	\item $10 \cdots \times 10^{-3}=10^{5}$;
	\item $10^{-4} \times 10 \cdots=10^{-3}$;
	\item $\dfrac{10^{2}}{10 \cdots}=10^{-2}$;
	\item $\dfrac{10^{-4}}{10^{\cdots}}=10^{4}$;
	\item $3 x \times \ldots=12 x^{3}$.
\end{enumMet}
\end{multicols}

\exercice{ $-$ ordre de grandeur}

Dans chaque cas, donner un ordre de grandeur du résultat.
\begin{multicols}{2}
	\begin{enumMet}
	\item $\mathrm{S}=4,98+8,01+9,96+26,02$;
	\item $\mathrm{S}=7,91+9,04+10,94+15,07$.
\end{enumMet}
\end{multicols}


\exercice{ $-$ équation de la forme $x^{2}=a$}

Résoudre dans \R les équations suivantes.
\begin{multicols}{4}
	\begin{enumMet}
	\item $x^{2}=4$;
	\item $x^{2}=5$;
	\item $2 x^{2}+3=0$;
	\item $x^{2}-7=0$.
\end{enumMet}
\end{multicols}


\exercice{ $-$ Développer, réduire et ordonner}

Dans ce qui suit, $P$ est une fonction polynôme définie par tout nombre réel $x$.

Développer et réduire chacune des expressions suivantes, puis ordonner les polynômes obtenus dans l'ordre des puissances décroissantes de $x$ ou de $t$ :

\begin{multicols}{2}
	\begin{enumMet}
	\item $P(x)=(x+1)(x+1)-(x+2) x$;
	\item $P(t)=(2 t-3)(2 t+3)-(2 t-3)^{2}$;
	\item $P(x)=(-x-1)^{2}+(x+1)^{2}$;
	\item $P(x)=x(3 x+2)^{2}-(3 x-2)^{2}$;
	\item $P(t)=\left(t-\dfrac{1}{2}\right)^{2}-\left(t-\dfrac{1}{2}\right)\left(t+\dfrac{1}{2}\right)$;
	\item $P(t)=\dfrac{t-5}{3}+\dfrac{2-t}{6}-t^{2}$.
\end{enumMet}
\end{multicols}

\exercice{ $-$ Factoriser, puis résoudre une équation.}

Mettre $P(x)$ sous forme d'un produit de facteurs du premier degré, puis résoudre l'équation $P(x)=0$.

\begin{multicols}{2}
	\begin{enumMet}
	\item $P(x)=x^{2}-4 x$.
	\item $P(x)=(x-3)(x-2)-(2 x-1)(x-3)$.
	\item $P(x)=9 x^{2}-6 x+1$.
	\item $P(x)=x^{2}-100$.
	\item $P(x)=x^{2}-5$.
	\item $P(t)=(7 t-3)^{2}-25$.
	\item $P(x)=(2 x+1)^{2}-(-x+3)^{2}$.
	\item $P(x)=(2 x+3)^{2}+(x-5)(2 x+3)$.
	\item $P(x)=(2 x-7)(-5 x+9)+4 x^{2}-49$.
\end{enumMet}
\end{multicols}

\exercice{ $-$ inéquations et tableau de signes}

Résoudre deux inéquations pour déterminer le signe d'une expression factorisée du second degré *

\begin{enumerate}
	\item Résoudre dans $\mathbb{R}$ l'inéquation $4 x+1 \geq 0$.
	\item Résoudre dans $\mathbb{R}$ l'inéquation $-2 x+3 \geq 0$.
	\item Compléter, après l'avoir reproduit, le tableau de signes suivant qui donne le signe $\operatorname{de}(4 x+1)(-2 x+3)$.
\end{enumerate}

\begin{center}
	\begin{tabular}{|l|rrrrrr|}
		\hline
		\multicolumn{1}{|c|}{$x$} & $-\infty$ & $-\frac{1}{4}$ &  & $\frac{3}{2}$ &  & $+\infty$ \\
		\hline
		Signe de $4 x+1$ & $?$ & 0 & $?$ &  & $?$ &  \\
		\hline
		Signe de $-2 x+3$ & $?$ &  & $?$ & 0 & $?$ &  \\
		\hline
		Signe de $(4 x+1)(-2 x+3)$ & $?$ & 0 & $?$ & 0 & $?$ &  \\
		\hline
	\end{tabular}
\end{center}

 
\end{document}

